$!A \ident (!B \lif !C)$ असेल, तर
$\pi$ चा शेवट असा असतो:
\[
\AxiomC{}
\RightLabel{$\pi_1'$}
\Deduce$!B, \Gamma \fCenter \Delta, !C$
\RightLabel{\RightR{\lif}}
\UnaryInf$\Gamma \fCenter \Delta, !B \lif !C$
\LeftSubproofLabel{$\pi_1$}
\AxiomC{}
\RightLabel{$\pi_2'$}
\Deduce$!B \lif !C, \Pi \fCenter \Lambda, !B$
\AxiomC{}
\RightLabel{$\pi_2''$}
\Deduce$!C, \Pi \fCenter \Lambda$
\RightLabel{\LeftR{\lif}}
\BinaryInf$!B \lif !C, \Pi \fCenter \Lambda$
\RightSubproofLabel{$\pi_2$}
\RightLabel{\Cut}
\BinaryInf$\Gamma, \Pi \fCenter \Delta, \Lambda$
\DisplayProof
\]
\Cut ने संपणारी !!{proof} पुढीलप्रमाणे:
\[
\AxiomC{}
\RightLabel{$\pi_1'$}
\Deduce$!B, \Gamma \fCenter \Delta, !C$
\RightLabel{\RightR{\lif}}
\UnaryInf$\Gamma \fCenter \Delta, !B \lif !C$
\LeftSubproofLabel{$\pi_1$}
\AxiomC{}
\RightLabel{$\pi_2'$}
\Deduce$!B \lif !C, \Pi \fCenter \Lambda, !B$
\RightLabel{\Cut}
\BinaryInf$\Gamma, \Pi \fCenter \Delta, \Lambda, !B$
\DisplayProof
\]
तिची कटची !!{height}
$\pi$ च्या कट-उंचीपेक्षा कमी आहे.
विगमन गृहीतकाने
$\Gamma, \Pi \fCenter \Delta, \Lambda, !B$
याची !!a{proof}~$\pi_1''$ मिळते.
\Cut ने संपणारी दुसरी
!!{proof} अशी आहे:
\[
\AxiomC{}
\RightLabel{$\pi_1'$}
\Deduce$!B, \Gamma \fCenter \Delta, !C$
\AxiomC{}
\RightLabel{$\pi_2''$}
\Deduce$!C, \Pi \fCenter \Lambda$
\RightLabel{\Cut}
\BinaryInf$!B, \Gamma, \Pi \fCenter \Delta, \Lambda$
\DisplayProof
\]
तिची कट-कोटी आणि
कटची !!{height} या दोन्ही
$\pi$ च्या कट-कोटी आणि
कटच्या !!{height} पेक्षा कमी आहेत.
म्हणून विगमन गृहीतक आणि
दुर्बलीकरण वापरून
$!B, \Gamma, \Pi, \Pi \fCenter \Delta, \Lambda,
\Lambda$ याची
~\Log{G3c} मधील !!{proof}~$\pi_2'''$
मिळते. आता पुढील
\Cut ने संपणाऱ्या
!!{proof} वर विगमन
गृहीतक लावा:
\[
\AxiomC{}
\RightLabel{$\pi_1''$}
\Deduce$\Gamma, \Pi \fCenter \Delta, \Lambda, !B$
\AxiomC{}
\RightLabel{$\pi_2'''$}
\Deduce$!B, \Gamma, \Pi, \Pi \fCenter \Delta, \Lambda,
\Lambda$
\RightLabel{\Cut}
\BinaryInf$\Gamma, \Gamma, \Pi, \Pi, \Pi \fCenter \Delta, \Delta, \Lambda, \Lambda, \Lambda$
\DisplayProof
\]
(हिची कट-कोटी
$< \cutr{\pi}$ आहे.)
त्यातून
$\Gamma,
\Gamma, \Pi, \Pi, \Pi \fCenter \Delta, \Delta, \Lambda, \Lambda, \Lambda$
याची !!a{proof} मिळते.
मग \olref[seq][inv]{prop:G3c-cont-adm}
लावून
$\Gamma, \Pi \fCenter \Delta, \Lambda$
याची !!{proof}~$\pi'$ मिळते.
